%!TEX root = thesis.tex
\chapter{Zusammenfassung}
\label{fazit und ausblick}

Große Sprachmodelle (LLMs) gewinnen zunehmend an Bedeutung bei der automatisierten Softwareentwicklung. Insbesondere die Erstellung von Quellcodes anhand unstrukturierter natürlichsprachlicher Anweisungen \textbf{"Vibe Coding"} eröffnet neue Möglichkeiten zur Unterstützung des Softwareentwicklungsprozesses. Ziel dieser Arbeit war die Untersuchung, welchen Einfluss unterschiedliche Prompt-Strategien und verschiedene LLMs auf die Qualität automatisch generierter Softwaresysteme haben.

Als Untersuchungsgrundlage diente das bereits entwickelte Softwaresystem \textit{Passat}, ein internes Buchungssystem zur Verwaltung von Gästen, Ressourcen und Buchungen auf einem Schiff. Für dieses System wurden das vorhandene Benutzerhandbuch sowie die funktionalen Anforderungen als Eingabe für die Sprachmodelle verwendet. Die Erstellung der Software erfolgte mit drei verschiedenen Prompt-Varianten – einem einfachen Basis-Prompt, einem strukturierten Prompt mit Kontextinformationen sowie einem strukturierten Prompt mit zusätzlicher Rollendefinition – und den drei Sprachmodellen GPT-5.2, Gemini 3.1 Pro und GPT-5.3-Codex. Daraus ergaben sich insgesamt neun Experimente.


Die Ergebnisse dieser Arbeit verdeutlichen, dass der erfolgreiche Einsatz großer Sprachmodelle in der Softwareentwicklung sowohl von der Leistungsfähigkeit des verwendeten Modells als auch von einer sorgfältigen Gestaltung der Prompts abhängt. Damit leistet die Arbeit einen Beitrag zum besseren Verständnis des Zusammenspiels von Prompt Engineering und LLMs bei der automatisierten Generierung von Softwaresystemen.

Eine wesentliche Herausforderung im Rahmen dieser Arbeit besteht in der nicht-deterministischen Funktionsweise großer Sprachmodelle, die sich während der Durchführung der Experimente wiederholt zeigte. Bei identischen Eingaben (Prompts) und demselben Sprachmodell können wiederholte Ausführungen unterschiedliche Ergebnisse liefern. Dadurch ist die Reproduzierbarkeit der generierten Software nur eingeschränkt gewährleistet.


Ein weiteres Problem bei der automatisierten Generierung von Softwaresystemen mit großen Sprachmodellen bestand in der Konsistenz zwischen den einzelnen Systemkomponenten. Werden Frontend, Backend und Datenbank unabhängig voneinander generiert, berücksichtigt das Sprachmodell die Schnittstellen zwischen den Komponenten nicht immer vollständig. Dadurch können Unterschiede in den Datenstrukturen, API-Endpunkten oder den verwendeten Attributnamen entstehen, was die Kommunikation zwischen den Komponenten beeinträchtigt.

Eine mögliche Vorgehensweise zur Reduzierung dieses Problems besteht darin, nicht das gesamte Softwaresystem gleichzeitig generieren zu lassen, sondern bereits vorhandene Komponenten als Grundlage zu verwenden. Beispielsweise kann ein bestehendes Backend einschließlich der Datenmodelle und REST-Schnittstellen dem Sprachmodell bereitgestellt werden, sodass anschließend nur das dazu passende Frontend generiert wird. Auf diese Weise kann sich das Modell an den vorhandenen Datenstrukturen und Schnittstellen orientieren, wodurch die Konsistenz zwischen Frontend und Backend verbessert und Integrationsfehler reduziert werden können.

Ein weiteres beobachtetes Problem betrifft die Generierung komplexer Prozesse, die aus mehreren aufeinander aufbauenden Schritten bestehen. Die Ergebnisse zeigen, dass große Sprachmodelle bei der Umsetzung solcher mehrstufigen Abläufe dazu neigen, einzelne Prozessschritte auszulassen, unvollständig zu implementieren oder in einer anderen Reihenfolge auszuführen. Mit zunehmender Komplexität des Entwicklungsprozesses steigt die Wahrscheinlichkeit, dass der generierte Ablauf nicht mehr vollständig den definierten Anforderungen entspricht.
Aufgrund der nur begrenzten Anzahl verfügbarer Output Tokens, konnte diese Vorgehensweise nur mit dem Modell GPT-5.3-Codex  vollständig untersucht werden. Eine umfassende Evaluation weiterer leistungsfähiger Sprachmodelle war daher nicht möglich.

Darüber hinaus wurde in einzelnen Experimenten beobachtet, dass die Generierung komplexer Prozesse vorzeitig beendet wurde oder keine vollständige Lösung erzeugt werden konnte. Für die Entwicklung komplexer Softwaresysteme empfiehlt es sich daher, umfangreiche Entwicklungsaufgaben in kleinere, klar definierte Teilschritte zu unterteilen und die generierten Ergebnisse nach jedem Schritt zu überprüfen. Auf diese Weise können Fehler frühzeitig erkannt und die Konsistenz des Gesamtsystems verbessert werden.
Außerdem wurde beobachtet, dass große Sprachmodelle bei umfangreichen Entwicklungsaufgaben mit zunehmender Anzahl an Interaktionen teilweise die Konsistenz bereits implementierter Komponenten verlieren. Änderungen an späteren Entwicklungsphasen können dazu führen, dass zuvor korrekt implementierte Funktionalitäten unbeabsichtigt verändert oder beeinträchtigt werden. Dadurch steigt der Aufwand für die Qualitätssicherung und Fehlerbehebung im weiteren Entwicklungsverlauf.
Aus wissenschaftlicher Sicht verdeutlicht der beobachtete Fehler, dass KI-generierte Software trotz syntaktisch korrekter Codegenerierung semantische Abhängigkeiten zwischen den einzelnen Softwarekomponenten nicht immer vollständig berücksichtigt. Dies kann dazu führen, dass Funktionalitäten erst zur Laufzeit fehlschlagen. Hier zeigen die vorliegenden Ergebnisse, dass die Bewertung KI-generierter Software nicht allein auf der Codequalität basieren sollte, sondern auch die Konsistenz zwischen Datenmodell, Geschäftslogik und Benutzeroberfläche berücksichtigen muss. Eine manuelle Validierung der erzeugten Softwareartefakte bleibt daher weiterhin erforderlich.
Zur Reduzierung dieses Problems bietet sich die Integration großer Sprachmodelle in eine integrierte Entwicklungsumgebung (IDE) an. Danach sollte nach jedem abgeschlossenen Entwicklungsschritt das Sprachmodell aufgefordert werden, geeignete automatisierte Tests für die neu implementierte Funktionalität zu erstellen. Anschließend sollten sowohl die neu erzeugten Tests als auch alle zuvor erstellten Tests ausgeführt werden. Schlägt einer der Tests fehl, sollten die identifizierten Fehler zunächst behoben und anschließend sämtliche Tests erneut ausgeführt werden. Erst wenn alle Tests erfolgreich bestanden werden, sollte mit dem nächsten Entwicklungsschritt fortgefahren werden.

Durch diese kontinuierliche Test- und Validierungsstrategie kann sichergestellt werden, dass neue Änderungen die bereits implementierten Funktionalitäten nicht unbeabsichtigt beeinträchtigen. Gleichzeitig ermöglicht dieses Vorgehen eine frühzeitige Erkennung von Integrationsfehlern und trägt zu einer höheren Qualität und Stabilität der generierten Software bei. \\

\section{Ausblick}
Die im Rahmen dieser Arbeit entwickelte Vorgehensweise erwies sich als effektiv für die schrittweise Generierung komplexer Softwaresysteme mit großen Sprachmodellen. Durch die iterative Entwicklung einzelner Komponenten sowie deren kontinuierliche Überprüfung konnte die Qualität der generierten Software deutlich verbessert werden. Gleichzeitig ist dieser Ansatz jedoch mit einem hohen Verbrauch von Output Tokens verbunden. Die wiederholte Interaktion mit den Sprachmodellen sowie die Generierung und Validierung einzelner Entwicklungsschritte führen zu einem erhöhten Bedarf an Tokens, Rechenleistung sowie Kosten und Energieverbrauch, wodurch insbesondere leistungsfähige kommerzielle Modelle erhebliche Nutzungskosten verursachen können.

Ein wesentlicher Aspekt, der in zukünftigen Arbeiten berücksichtigt werden sollte, besteht darin, die entwickelte Methodik auf neuere und leistungsfähigere Sprachmodelle, beispielsweise Claude Sonnet oder Claude Opus, zu übertragen und deren Ergebnisse systematisch mit den in dieser Arbeit untersuchten Modellen zu vergleichen. Dabei könnte untersucht werden, inwieweit moderne Modelle bereits mit einfachen oder strukturierten Basis-Prompts qualitativ vergleichbare oder bessere Ergebnisse erzielen als die in dieser Arbeit eingesetzten Modelle. Ein solcher Vergleich würde zusätzliche Erkenntnisse über den Einfluss der Modellgeneration auf die Qualität der automatisierten Softwareentwicklung liefern.

Ein weiterer Forschungsansatz ergibt sich aus der Kombination von Test-Driven Development (TDD) und großen Sprachmodellen. Anstatt das Sprachmodell ausschließlich auf Grundlage der funktionalen Anforderungen zur Generierung einer Software zu verwenden, könnte zunächst eine vollständige Testsuite erstellt werden. Anschließend würden sowohl die Anforderungen als auch die Testfälle dem Sprachmodell bereitgestellt mit dem Ziel, ein Softwaresystem zu generieren, das sämtliche Tests erfolgreich besteht und eine definierte Testabdeckung, beispielsweise von mindestens 80\,\%, erreicht. Ein solcher Ansatz könnte die Qualität und Zuverlässigkeit der generierten Software weiter verbessern, da die Implementierung unmittelbar anhand objektiver Testkriterien validiert wird. Gleichzeitig ergeben sich daraus mehrere Forschungsfragen, die in zukünftigen Arbeiten untersucht werden können.

Hierzu zählt insbesondere die Frage, ob für verschiedene Entwicklungskriterien, Frameworks oder Anwendungsdomänen jeweils spezifische Testsuiten entwickelt werden müssen oder ob sich ein allgemeingültiger Testansatz für unterschiedliche Softwaresysteme realisieren lässt.

Weiterhin stellt sich die Frage, ob die Testfälle von Softwareentwicklern manuell erstellt oder ebenfalls durch große Sprachmodelle generiert werden sollten. Beide Vorgehensweisen unterscheiden sich hinsichtlich Entwicklungsaufwand, Kosten, Zeitbedarf und Qualität der erzeugten Tests. Eine systematische Untersuchung dieser Aspekte könnte wertvolle Erkenntnisse darüber liefern, welche Strategie für den praktischen Einsatz großer Sprachmodelle in der Softwareentwicklung sowohl wirtschaftlich als auch technisch am sinnvollsten ist.
